Recall from the setting (notation section of the paper; Blueprint definition ``setting) that $G_n=\langle\sigma\rangle$ is cyclic of order $2^n$ and $\tau=\sigma^{2^{n-1}}$ is its unique element of order $2$; $B_{n-1}$ is the fixed field of $\tau$, so $\operatorname{Gal}(B_n/B_{n-1})=\{1,\tau\}$ and the relative norm is $N_{n/n-1}(x)=x\cdot\tau(x)$ for $x\in B_n$. By definition $RE^+_n=\{\varepsilon\in E_n:\ N_{n/n-1}\varepsilon=1\}$, and $E_{n-1}=E_n\cap B_{n-1}$ is the set of units fixed by $\tau$.

Let $u\in RE^+_n\cap E_{n-1}$. Because $u\in E_{n-1}$ we have $\tau(u)=u$; because $u\in RE^+_n$ we have $u\cdot\tau(u)=1$. Substituting the first relation into the second gives $u\cdot u=u^2=1$. (This is \path{bridge_fixed_norm_sq}, stated for an arbitrary abelian group $G$ with a homomorphism $\tau:G\to G$: from $\tau u=u$ and $u\,\tau u=1$ one obtains $u^2=1$; it is applied with $G=E_n$.)

In the field $B_n$ the identity $u^2=1$ reads $(u-1)(u+1)=0$, and a field has no zero divisors, so $u=1$ or $u=-1$ (\path{torsion_pm_one}; only the field axioms of $B_n$ are used, not that $B_n$ is totally real). Hence $RE^+_n\cap E_{n-1}\subseteq\{\pm1\}$. Conversely $\pm1$ are units of $B_{n-1}$ with $N_{n/n-1}(\pm1)=(\pm1)^2=1$, so in fact $RE^+_n\cap E_{n-1}=\{\pm1\}$; only the inclusion is used later.
